\section{結果：北海道}
\label{sec:hokkaido}

対象は \ang{40.0}--\ang{46.5}N / \ang{138.0}--\ang{149.0}E、
\ang{0.375} 格子で 570 点。周辺の日本海・オホーツク海・太平洋を含む。
代表地点は \texttt{sites.json} に登録した \textbf{VOISING}
（\ang{42.1823}N, \ang{142.8510}E、日高・浦河の東）である。
最寄り格子までの距離は \SI{4.7}{\kilo\metre}。
どの地点を代表にするかはコード側で名前指定してあり
（\texttt{csv\_out.REFERENCE\_SITE}）、
\texttt{sites.json} の並び順には依存しない。

\subsection{まず標本数を確認する}

\begin{figure}[H]
\centering
\includegraphics[width=0.95\linewidth]{hokkaido/F0_sample_counts.pdf}
\caption[標本数]{各ビンに入った解析値の数。前半は必ず 15 日ぶんで 600 個、後半は月の長さで変わり 1 月は 640 個、2 月は閏年が 3 回あるので 532 個になる。この図が合わなければ取得か集計が壊れているので、他のどの図よりも先に確認する。}
\label{fig:f0-hokkaido}
\end{figure}


図 \ref{fig:f0-hokkaido} の標本数は設計どおりである。
前半が各 600 個、後半は月の長さに応じて 532--640 個。
一方、有効標本数（第 \ref{sec:neff} 節）の中央値は \textbf{77.2} で、
生の標本数の 13 \% にすぎない。
つまり「10 年分の 600 観測」は、独立な天気としては 77 個程度の情報しか持っていない。
以降の平均値は、この前提で読む必要がある。

\subsection{年間を通した姿}

\begin{figure}[H]
\centering
\includegraphics[width=1.0\linewidth]{hokkaido/F1_alt_half_speed.pdf}
\caption[時期×高度の風速と風向]{上：平均風速。下：平均風向（吹いてくる向き）。風向には円環用のカラーマップを使っている。順次カラーマップを使うと \ang{0} と \ang{360} の境目に実在しない断層が現れるためである。灰色は定常度が 0.3 未満で塗りを外した領域で、そのビンではそもそも風向が定まっていない。}
\label{fig:f1-hokkaido}
\end{figure}

\begin{figure}[H]
\centering
\includegraphics[width=1.0\linewidth]{hokkaido/F3_profiles_by_month.pdf}
\caption[月別の高度プロファイル]{月ごとの風速プロファイル。線は各ビンの中央値、灰色の帯はその月の 2 ビンを通した p10--p90 の広がり。帯の幅が平均からのずれの大きさを表しており、線だけを見て計画を立ててはいけない。}
\label{fig:f3-hokkaido}
\end{figure}


対流圏界面付近（\SI{200}{\hecto\pascal}、約 \SI{11.8}{\kilo\metre}）の平均風速は
冬 \SI{40.2}{\metre\per\second}、春 \SI{35.0}{\metre\per\second}、
夏 \SI{29.9}{\metre\per\second}、\textbf{秋 \SI{44.5}{\metre\per\second}} である。

冬ではなく秋が最も強いのは、緯度による。
東アジアのジェットは冬に最も強くなるが、同時に南へ下がる。
北海道の緯度帯には、ジェットが南下しきる前の 10--11 月に核が通過する。
後述する和歌山（冬に \SI{75.5}{\metre\per\second}）と比べると、この違いははっきりする。

定常度は対流圏では総じて高い（\SI{200}{\hecto\pascal} で 0.81--0.96）。
つまり\textbf{北海道上空の対流圏の風向は年間を通して安定している}。
一方で、上層には風向が定まらない時期がある。

\subsection{上層では風向が定まらない時期がある}

\begin{figure}[H]
\centering
\includegraphics[width=1.0\linewidth]{hokkaido/F2_alt_half_constancy.pdf}
\caption[定常度]{定常度（ベクトル平均風速 ÷ スカラー平均風速）。1 に近いほど風向が安定し、0 に近いほどビンのなかで風向が入れ替わっている。淡い帯が成層圏の風の反転している高度で、季節とともに上下する。}
\label{fig:f2-hokkaido}
\end{figure}


\SI{10}{\hecto\pascal}（約 \SI{31}{\kilo\metre}、破裂高度帯）の定常度は
夏に 0.99 と極めて高い一方、春は 0.36、秋は 0.31 しかない。
特に定常度が低いのは 10 月前半（0.15）、5 月前半（0.17）、
11 月後半（0.21）、10 月後半（0.23）、3 月前半（0.26）である。

これらは成層圏の風が東西逆転する時期にあたる。
そのビンの 10 年ぶんの解析には「強い西風の日」と「強い東風の日」が混ざっており、
平均するとほぼ打ち消し合う。
図 \ref{fig:f1-hokkaido} の下パネルでこれらのビンが灰色になっているのはそのためで、
\textbf{データが無いのではなく、平均風向という量に意味がない}。

\begin{figure}[H]
\centering
\includegraphics[width=1.0\linewidth]{hokkaido/F7-1_wind_roses.pdf}
\caption[風配図（前半）]{代表格子点の風配図、\textbf{月の前半}（1--15 日）。
平均と楕円では見せられない「分布そのもの」を示す。
ローブが 2 つに割れている場面が、平均風がほぼゼロで分散だけが大きくなる時期にあたる。}
\label{fig:f7-hokkaido-h1}
\end{figure}

\begin{figure}[H]
\centering
\includegraphics[width=1.0\linewidth]{hokkaido/F7-2_wind_roses.pdf}
\caption[風配図（後半）]{同じ格子点の\textbf{月の後半}（16 日--月末）。
前半（図\ref{fig:f7-hokkaido-h1}）と並べると、\textbf{2 枚はよく似ている}。
\SI{30}{\hecto\pascal} の卓越風向は 1・4・7・10 月のいずれも前半と後半で同じである。
これは第 \ref{sec:validation} 章の結論と整合する。
月をさらに細かく割っても、新しい構造は出てこない。}
\label{fig:f7-hokkaido-h2}
\end{figure}

風配図（図 \ref{fig:f7-hokkaido-h1}、\ref{fig:f7-hokkaido-h2}）はこの状況を直接見せる。
平均と分散楕円では 2 つの塊を 1 つの楕円に潰してしまうが、
風配図なら塊が 2 つあればそう見える。

\subsection{冬の 10 hPa が東風であること}

一つ、事前の予想と食い違った点を記しておく。

帯状平均で見れば、中緯度の成層圏は冬に西風（極夜ジェット）である。
ところが本統計では、北海道上空の \SI{10}{\hecto\pascal} は
1 月に東風（\SI{-3.8}{\metre\per\second}）になる。
2016、2017、2018 の各年を個別に見ても、いずれも弱い東風である。
さらに同じ高度を経度方向に見ると、
東へ行くほど東風が強まる勾配がある。

これは誤りではなく、この経度帯の実際の気候である。
真冬の成層圏には北太平洋上に高気圧性の循環（アリューシャン高気圧）が停滞し、
その西縁にあたる日本の東側では帯状平均と逆向きの風になる。
\textbf{帯状平均の知識を、一部の経度だけを切り取った領域にそのまま当てはめてはいけない}という例であり、
このパイプラインの検証項目もその点で修正した。

\subsection{領域の中で答えが変わるか}

\begin{figure}[H]
\centering
\includegraphics[width=1.0\linewidth]{hokkaido/F6_spatial_spread.pdf}
\caption[領域内のばらつき]{領域内での平均風速の幅（最大 $-$ 最小）。これが小さい高度・時期では、領域のどこで上げても風はほぼ同じで、代表 1 地点のプロファイルで語ってよい。大きいところではそうではない。}
\label{fig:f6-hokkaido}
\end{figure}


領域内での平均風速の幅（最大 $-$ 最小）は最大で \SI{31.7}{\metre\per\second} に達する。
これは無視できない大きさで、
\textbf{北海道の領域は「代表 1 地点で語れる」ほど一様ではない}。
南端（\ang{40}N、東北北部）と北端（\ang{46.5}N、宗谷沖）では、
特に冬のジェットの強さが大きく違う。

したがって VOISING の数値を道北や道東にそのまま適用してはいけない。
格子ごとの値は NetCDF と Parquet にあるので、必要な地点を引くこと。

\subsection{地図で見る}

\begin{figure}[H]
\centering
\includegraphics[width=1.0\linewidth]{hokkaido/M1-200_month_grid_200hPa.pdf}
\caption[200 hPa の月別分布図]{\SI{200}{\hecto\pascal}（約 \SI{12}{\kilo\metre}）の平均風を月ごとに並べたもの。塗りが風速、矢印が風向。ジェットが季節とともに動く様子が地図の上で見える。}
\label{fig:m1200-hokkaido}
\end{figure}

\begin{figure}[H]
\centering
\includegraphics[width=1.0\linewidth]{hokkaido/M1-030_month_grid_30hPa.pdf}
\caption[30 hPa の月別分布図]{同じく \SI{30}{\hecto\pascal}（約 \SI{24}{\kilo\metre}）。対流圏とは逆の季節変化をしており、夏に東風が強まる。}
\label{fig:m1030-hokkaido}
\end{figure}

\begin{figure}[H]
\centering
\includegraphics[width=1.0\linewidth]{hokkaido/M3-200_direction_map_200hPa.pdf}
\caption[風向の分布図]{\SI{200}{\hecto\pascal} の平均風向。円環用のカラーマップを使い、定常度 0.3 未満は灰色にしている。灰色の領域では平均風向に意味がない。}
\label{fig:m3200-hokkaido}
\end{figure}

\begin{figure}[H]
\centering
\includegraphics[width=1.0\linewidth]{hokkaido/M2_constancy_map.pdf}
\caption[定常度の分布図]{高度と季節ごとの定常度。濃いところでは風向が安定していて計画が立てやすく、淡いところではビンのなかで風向が入れ替わる。}
\label{fig:m2-hokkaido}
\end{figure}

\begin{figure}[H]
\centering
\includegraphics[width=0.92\linewidth]{hokkaido/M4-01_wind_ellipse_map.pdf}
\caption[風のばらつきを地図に]{1 月前半、\SI{200}{\hecto\pascal} の平均風ベクトルと、その先端に描いた $1\sigma$ の散らばり楕円。丸い楕円はどの向きにもばらつくこと、細長い楕円は主に強さがばらつくことを意味する。}
\label{fig:m401-hokkaido}
\end{figure}


\subsection{落下位置}

\begin{figure}[H]
\centering
\includegraphics[width=1.0\linewidth]{hokkaido/M5_drift_month_grid.pdf}
\caption[月別の平均落下位置]{各格子点から上げた場合の平均落下位置。\textbf{矢印は実寸}で、矢先がそのまま予測落下点である。色が平均の飛行距離。全パネルが同じ枠と同じ色スケールなので、月どうしを直接比べられる。}
\label{fig:m5-hokkaido}
\end{figure}

\begin{figure}[H]
\centering
\includegraphics[width=0.92\linewidth]{hokkaido/M7-01_drift_map.pdf}
\caption[1 月の落下位置]{1 月前半の平均落下位置を拡大したもの。枠は集計対象の領域で、矢先はその外へ出ていることが多い。}
\label{fig:m701-hokkaido}
\end{figure}

\begin{figure}[H]
\centering
\includegraphics[width=0.92\linewidth]{hokkaido/M7-07_drift_map.pdf}
\caption[7 月の落下位置]{同じく 7 月前半。1 月と比べると距離も向きも大きく違う。}
\label{fig:m707-hokkaido}
\end{figure}

\begin{figure}[H]
\centering
\includegraphics[width=1.0\linewidth]{hokkaido/M6_drift_percentile_map.pdf}
\caption[荒れた日の飛行距離]{飛行距離の 90 パーセンタイル。平均 \SI{150}{\kilo\metre} で p90 が \SI{400}{\kilo\metre} という状況と、平均 \SI{150}{\kilo\metre} だけを見て計画することは、運用上まったく違う。}
\label{fig:m6-hokkaido}
\end{figure}

\begin{figure}[H]
\centering
\includegraphics[width=0.92\linewidth]{hokkaido/M8_landing_ellipse_map.pdf}
\caption[代表地点からの落下分布]{1 地点から上げた場合の落下位置を季節別に重ねたもの。楕円は 50 / 90 / 95 \% の範囲で、各解析時刻ごとに求めた着地点の共分散から描いている。散らばりが 2 変量正規に従うという仮定は残る。灰色が陸、白が海である。}
\label{fig:m8-hokkaido}
\end{figure}


飛行モデル（第 \ref{sec:drift} 節、上昇 \SI{5}{\metre\per\second}、
破裂 \SI{30}{\kilo\metre}、降下 \SI{6}{\metre\per\second}）での飛行時間は約 \num{138} 分。
領域平均の飛行距離は年間平均 \SI{143}{\kilo\metre} で、季節差が非常に大きい。

\begin{table}[H]
\centering
\caption{北海道：飛行距離が最短・最長のビン（領域平均）。
方位は機体が\textbf{進む}向きで、風向とは \ang{180} 逆である。}
\begin{tabular}{lrrr}
\toprule
時期 & 平均距離 [\si{\kilo\metre}] & 方位 [度] & p90 [\si{\kilo\metre}] \\
\midrule
7 月 後半 & 62 & 123（東南東） & 102 \\
7 月 前半 & 67 & 102（東南東） & 108 \\
8 月 前半 & 72 & 96（東） & 122 \\
\midrule
12 月 前半 & 215 & 75（東北東） & 292 \\
11 月 後半 & 205 & 78（東北東） & 284 \\
12 月 後半 & 192 & 82（東） & 269 \\
\bottomrule
\end{tabular}
\end{table}

この差は運用の判断に直接つながる。

\begin{itemize}
  \item \textbf{7 月が最も扱いやすい。} 飛行距離が冬の 4 分の 1 になり、
        p90 でも \SI{100}{\kilo\metre} 程度に収まる。
  \item \textbf{11 月から 12 月は避けたほうがよい。} 平均で \SI{200}{\kilo\metre} を超え、
        p90 は \SI{300}{\kilo\metre} に達する。方位は東北東で一定なので、
        十勝から上げた機体はほぼ確実に太平洋に落ちる。
  \item 図 \ref{fig:m8-hokkaido} のとおり、VOISING から上げた場合、
        \textbf{平均落下点はどの季節でも海上}である。
        VOISING は日高の南岸で太平洋に面しており、
        機体は東へ運ばれるので、そのまま沖へ出る。
\end{itemize}

季節で変わるのは「海に落ちるかどうか」ではなく「岸からどれだけ離れるか」である。

\begin{table}[H]
\centering
\caption{VOISING から上げた場合の平均落下点と、そこから最寄りの海岸までの距離。
どの時期も海上だが、夏だけは岸のすぐそばに落ちる。}
\label{tab:voising-offshore}
\sisetup{table-format=3.0}
\begin{tabular}{@{}llSS@{}}
\toprule
月 & ビン & {飛行距離 [\si{\kilo\metre}]} & {岸まで [\si{\kilo\metre}]} \\
\midrule
1 月  & 前半 & 187 & 80 \\
1 月  & 後半 & 184 & 75 \\
4 月  & 前半 & 156 & 67 \\
4 月  & 後半 & 155 & 74 \\
7 月  & 前半 & 71  & 4 \\
7 月  & 後半 & 62  & 2 \\
10 月 & 前半 & 184 & 59 \\
10 月 & 後半 & 185 & 55 \\
\bottomrule
\end{tabular}
\end{table}

\textbf{7 月だけは状況が大きく違う。} 平均落下点が岸から \SIrange{2}{4}{\kilo\metre} しかないので、
散らばり（図 \ref{fig:m8-hokkaido} の楕円）の陸側の裾は陸にかかる。
回収を陸上で行いたいなら、実質的に 7 月しか選択肢がない。
他の季節は \SIrange{55}{80}{\kilo\metre} 沖なので、船を用意するか、回収を諦めるかを決めておく必要がある。

\begin{whybox}{打ち上げ地点を変えるとこの結論は変わる}
これは VOISING という\textbf{特定の地点}の性質であって、北海道の性質ではない。

同じ北海道でも、内陸寄りの地点から上げれば同じ風でも陸上に落ちる。
第 \ref{sec:hokkaido} 章の他の図は領域全体の格子点を扱っているので
地点によらないが、この図と表だけは地点に依存する。

\texttt{out/csv/region\_sites\_hokkaido.csv} に領域内の全 114 地点と
最寄り格子までの距離を出してあるので、別の地点で見たい場合は
\texttt{csv\_out.REFERENCE\_SITE} を変えて再実行すればよい。
\end{whybox}

なお、ここでの散らばりは\textbf{風だけに由来する}。
気球が予定より早く割れたり粘ったりする効果は含まれていないので、
実際の散らばりはこれより大きい。

\subsection{日変化と鉛直シアー}

\begin{figure}[H]
\centering
\includegraphics[width=0.75\linewidth]{hokkaido/F8_diurnal.pdf}
\caption[日変化]{夜（03 JST）と昼（15 JST）の平均風速の差を、季節ごとに高度に対して描いたもの
（\SI{300}{\hecto\pascal} 以下の層のみ、領域平均）。横軸の目盛りに注意してほしい。
差は 3 領域・4 季節のどれでも \SI{1}{\metre\per\second} 未満で、
日々の変動（数十 \si{\metre\per\second}）に比べて小さい。
これが\textbf{時刻を区別せず平均してよい根拠}である。
ビンを時刻で 4 分割すると標本数も 4 分の 1 になるので、
必要のないところで分割しない理由がここにある。}
\label{fig:f8-hokkaido}
\end{figure}

\begin{figure}[H]
\centering
\includegraphics[width=1.0\linewidth]{hokkaido/F9_shear.pdf}
\caption[鉛直シアー]{隣り合う気圧面のあいだの風速差を高度差で割ったもの。
機体そのものはシアーで壊れないが、大きいところでは
高度の推定が少しずれるだけで水平の予測が大きく動く。
最大はジェットの上下と地表付近に出る。}
\label{fig:f9-hokkaido}
\end{figure}


夜（03 JST）と昼（15 JST）の平均風速の差は、最大でも \SI{0.9}{\metre\per\second} である
（図\ref{fig:f8-hokkaido}）。北海道では地表付近より上空（高度 \SI{9}{\kilo\metre} 付近、昼のほうが強い）
のほうが差が大きいが、それでも日々の変動に比べれば小さい。
第 \ref{sec:method} 章で上層の統計を解析時刻で分けなかったのは、
この図が根拠である。

鉛直シアーの最大は \SI{6.7}{(\metre\per\second)\per\kilo\metre}
（5 月前半、高度 \SI{0.1}{\kilo\metre}）。
地表付近を除くと \SI{5.4}{(\metre\per\second)\per\kilo\metre}
（10 月前半、高度 \SI{16.4}{\kilo\metre}）で、対流圏界面のすぐ上にあたる。

\subsection{既存シミュレータの ISA 仮定}

\begin{figure}[H]
\centering
\includegraphics[width=1.0\linewidth]{hokkaido/F10_isa_check.pdf}
\caption[ISA との比較]{気圧から高度への換算。左は JRA-3Q の実際の気圧面高度と、既存シミュレータが使う国際標準大気 ISA。右はその差。ISA は季節も緯度も天気も知らないので、季節によって系統的にずれる。}
\label{fig:f10-hokkaido}
\end{figure}


既存シミュレータ（\texttt{backend/bfps/core/atmosphere.py}）は
気圧から高度への換算に国際標準大気を使っている。
北海道での実際とのずれは最大 \SI{1036}{\metre}（夏）、冬でも \SI{520}{\metre} ある。

これは無視できる量ではない。
たとえば「高度 \SI{30}{\kilo\metre} で破裂」と設定した計算が、
夏には実際の \SI{29}{\kilo\metre} 相当を指してしまう。
高度そのものを扱う場面では、ISA ではなく本統計の
\texttt{alt\_geom\_mean\_m} を使うべきである。
